\documentclass{article}
\usepackage{amsmath}

\title{5. Főcsoport: A Nitrogén csoport}
\author{Haydu Lőrinc, Lengyel Zoltán Péter, Rovenszky Robin}
\date{Legutolsó frissítés: 2026. Március 13.}
\newtheorem{definition}{Definition}[section]

\begin{document}

\maketitle

\section{A Nitrogén csoport}
\begin{enumerate}
    \item Nitrogén (N)
    \item Foszfor (P)
    \item Arzén (As)
    \item Antimon (Sb)
    \item Bizmut (Bi)
\end{enumerate}

5 db vegyérték elektron $\implies$ a nemesgáz elektronszerkezetének eléréséhez:
\begin{itemize}
    \item 3 db kovalens kötés\\
    Pl.: Ammónia 
    % insert ammónia ábra here
    \item 3-szoros negatív töltésű anion (pl. $N^{3-}$ - nitrid-ion). \\
    Egyedi, mert lényegében csak a nátrium tud ennyi iont felvenni, mivel nagyon nagy az elektronegativitása.
    \item Az is lehetséges, hogy a nitrogén 5 db kovalens kötést létesítsen. Pl.: $HNO_3$ - salétromsav
    % insert salétromsav ábra here
    \\
    A nitrogén túllépett a nemesgáz elektroszerkezetén, de ez így is energetikailag kedvező.
\end{itemize}

\section{Nitrogén}
\subsection{Atomi jellemzők}
\begin{itemize}
    \item nagy EN (3.0)
\end{itemize}
\subsection{Molekulaszerkezet}
%insert Nitrogen atom here
Ez egy erős kovalens kötés $\rightarrow$ közönséges körülmények között nem reakcióképes
$\rightarrow$ apoláris molekula $\rightarrow$ diszperzioós kölcsönhatás $\rightarrow$ molekularácsos szilárd halmazállapotban
\subsection{Fizikai tulajdonságok}
\begin{itemize}
    \item alacsony olvadáspont és forráspont (-210C és -180C)
    \item vízben rosszul oldódik
\end{itemize}

\begin{itemize}
    \item (keszonbetegség $\rightarrow$ gázok oldhatósága nagyobb nyomás alatt jobb!)
\end{itemize}

\subsection{Kémiai tulajdonságok}
\begin{itemize}
    \item $N_2 + O_2 \rightarrow^{3000C}2NO$
    \item $\mathrm{N_2 + 3H_2 \ \overset{\text{Fe-katalizátor}}{\rightleftharpoons}\ 2NH_3 }$
\end{itemize}

\subsection{Biológiai kitekintő: Nitrogén körforgása a környezetben}
\begin{itemize}
    \item A nitrogén hármas kötése miatt majdnem semmi sem veszi fel, viszont nitrifikáló baktériumok fel tudják venni, ezek pedig szimbiózisban élnek bizonyos növényekkel (pl. bab, borsó, ill. egyéb hüvelyesek).
    \item Másik mód amivel a körforgásba kerül a $3000^\circ\text{C}$-n bekövetkező reakció (villámlás).
\end{itemize}

\subsection{Szervetlen vegyipar legfontosabb folyamatai}
\begin{itemize}
    \item Hidrogén-klorid gyártás
    \item Kénsac-gyártás
    \item Ammónia-gyártás
\end{itemize}


\subsection{Felhasználás}
\begin{enumerate}
    \item Ammónia gyártás
    \item Védőgáz (csomagolás, hegesztés)
    \item Hűtés (folyékony nitrogénnel)

\end{enumerate}
\subsection{Előállítás} 
A levegő cseppfolyósításával és frakcionális desztillációjával elkülönítikd

\begin{definition}[Frakcionált desztilláció.]
A forráspontkülönbségen alapuló elválasztási művelet.
\end{definition}

\subsection{A vegyületei}

\begin{itemize}
    \item Oxidjai
    \begin{enumerate}
        \item $N_2O$ (dinitrogén-monoxid)\\
        Színtelen, édeskés szagú gáz \begin{itemize}
            \item Felhasználása: \begin{itemize}
                \item Tejszínhab készítése (szifon)
                \item Érzéstelenítés orvoslás tudományban (szülési fájdalmak csökkentése)
                \item Méhészet (elbódítja a méheket mielőtt beavatkozik)
                \item Autóversenyek (nitro)
            \end{itemize}
            \item Hétköznapi neve: kéjgáz (nevetőgáz)
        \end{itemize}

        \item $NO$ (nitrogén-monoxid)
        \[
        N_2 + O_2 \xrightarrow{3000^\circ\mathrm{C}} 2NO
        \]
        Tovább alakul:
        \[
        2NO + O_2 \rightarrow 2NO_2
        \]
        Színtelen, szúrós szagú, mérgező gáz, az érfalak rugalmasságát szabályozza az emberi szervezetben $\implies$ a vérnyomásra hat (nitroglicerin tabletta)
        % insert nitroglicerin ábra here

        \item $NO_2$ (nitrogén-dioxid)\\
        % insert nitrogén-dioxid ábra here
        Vörösbarna, szúrós szagú gáz, erősen mérgező. A szennyezett városi legevő egyik mérgező összetevője.
        
    \end{enumerate}

    \item Oxosavai
    \begin{enumerate}
        \item Salétromossav
        \[
        2NO_2 + H_2O \rightarrow \underset{\text{salétromossav}}{HNO_2} + \underset{\text{salétromsav}}{HNO_3}
        \]
        Gyenge, bomlékony, egyértékű sav. Csak vizes oldatban létezik. Sói: nitritek, pl.: $KNO_2$ $\leftarrow$ kálium-nitrit $\leftarrow$ pác\textbf{só} $\leftarrow$ tartósítószer
    \end{enumerate}
\end{itemize}

\end{document}